\chapter{Introducción}

El propósito de este libro de investigación es analizar en profundidad los modelos de \gls{iot} abiertos y comprender cómo pueden proporcionar soluciones efectivas para la lectura y transmisión de datos que los objetos ofrecen a través de los sensores. Al comprender estos modelos, será posible controlar o leer parámetros de los objetos para aplicaciones inteligentes. Es importante destacar que el uso de sistemas de hardware y software abiertos ofrece numerosas ventajas, como la interoperabilidad entre diferentes fabricantes y el acceso a las arquitecturas y códigos necesarios para ajustar los proyectos de \gls{iot} a las necesidades específicas. Por lo tanto, este libro tiene como objetivo brindar a los lectores una comprensión clara de los modelos de IoT abiertos y su aplicación en diferentes campos.

Este libro es el resultado de un proyecto de investigación en el que se ha desplegado una red \gls{lora}. El objetivo del proyecto es desarrollar un sistema de nodos, concentradores y servidores, junto con software libre diseñado para la gestión de la red y el tratamiento de datos. El objetivo final es proporcionar resultados eficaces, desde la recolección de los datos hasta la entrega de la información a través de una interfaz web.

El libro está estructurado en capítulos que abordan temas específicos de la red \gls{lora}. En los capítulos 2 y 3,  se describen las generalidades y los los fundamentos de la red \gls{lora}, desde los protocolos hasta la descripción de cada capa de la red. Se comienza con la capa física, seguida de la capa de MAC, que se despliega en tres clases: Clase A, Clase B y Clase C. Por último, se describe la capa de aplicación. Se proporcionará una descripción detallada de las tramas de recepción y transmisión que involucran cada una de las capas de la red \gls{lora}. También se incluirá una descripción de las bandas de frecuencia asignadas a este estándar, así como los tipos de modulación y demodulación asignados a la transmisión y recepción. La característica más relevante de la red \gls{lora} es la seguridad de los datos, que es vital en el despliegue radioeléctrico de la información.

En el capítulo 4, se abordarán las arquitecturas de los nodos propuestos por el estándar \gls{lora}. Se describirán detalladamente cada uno de los bloques que componen el nodo. En primer lugar, se analizará la fuente de alimentación y se especificarán los márgenes de tolerancia requeridos para su correcto funcionamiento. A continuación, se estudiará la inteligencia del nodo a través de un microcontrolador que deberá cumplir con ciertos criterios para ser utilizado en el proyecto, como la capacidad de adquirir datos, el reloj, la gestión de los protocolos del nodo, los periféricos para la conexión de los sensores y la capacidad de procesamiento de datos, entre otros.

También se abordará el bloque de radio \gls{lora}, que se caracteriza por su bajo consumo de energía y se mostrará su arquitectura y el manejo de los periféricos. En este capítulo se destacará el diseño del nodo para el proyecto, que puede implementarse basado en un chipset \gls{lora}, un RF-MCU o un módem \gls{lora} interno o externo. Asimismo, se describirá la interfaz digital del nodo, la cual debe ser compatible con el estándar \gls{lora} y utilizar las principales interfaces como $I^2C$, SPI y UART. Se presentará una descripción detallada de la interfaz utilizada en el proyecto, así como la interfaz de hardware con la disposición sugerida para su correcto funcionamiento.

En el capítulo 5, además de introducir los sensores, se profundizará en la selección de los mismos para el proyecto en cuestión. Se explicará detalladamente cada uno de los sensores necesarios y su respectiva interfaz, desde los sensores de interfaz analógica, que miden señales analógicas como la temperatura o la humedad, hasta los sensores de interfaz digital que permiten la comunicación directa con el microcontrolador a través de protocolos como $I^2C$, SPI, entre otros.

En el caso específico del sensor de CO2, se describirán sus características y su interfaz de despliegue, así como también se hablará de la forma en que se medirá su concentración en el ambiente y cómo se utilizarán los datos recolectados para las aplicaciones inteligentes. 

Además, se abordará el tema de los sensores de geolocalización, explicando su funcionamiento y su importancia en el proyecto, así como también se describirán otros sensores utilizados en el proyecto, como pueden ser los sensores de luminosidad, los sensores de temperatura, los sensores de humedad, entre otros. En general, se explicará el uso de cada uno de los sensores y se describirán detalladamente sus características técnicas y su interfaz de conexión.

El capítulo 6 se enfocará en el concentrador LoRa, donde se describirán detalladamente los componentes que lo conforman, como el host, el núcleo y las funcionalidades básicas para su correcto funcionamiento. Como host se utilizará el minicomponente Raspberry Pi 3, que cuenta con una interfaz GPIO para la conexión de periféricos externos. La pasarela utilizada es la encargada de la interfaz de radio y se encarga de comunicarse con miles de sensores y a través del host, llevar esos datos a un servidor LoRaWAN para la gestión de los datos.

Se empleará el sistema operativo Alpine Linux, debido a que es muy liviano y está orientado a la seguridad, además de ser de propósito general. Se presentará una guía ilustrada de cómo configurar el sistema operativo para optimizar su desempeño. Se utilizará el software LoRa Packet Forwarder para enviar los paquetes de datos al servidor, y el software LoRa Server será el encargado de gestionar las tramas de datos, tanto descendentes como ascendentes, y en general, de la programación de la transmisión de los datos de enlace.

En cuanto al almacenamiento y respaldo de datos, se utilizará el software InfluxDB como base de datos. Por último, para la visualización de los datos, se empleará el software Grafana, el cual permite a los usuarios visualizar los datos en línea y en tiempo real. Todo esto se explicará detalladamente en el capítulo 4, incluyendo el proceso de instalación y configuración de cada uno de los componentes, así como la interacción entre ellos para asegurar el correcto funcionamiento del concentrador LoRa.

En el capítulo 7, se presenta la implementación de la red \gls{lwan} en su totalidad, una vez que se han realizado las pruebas individuales de cada uno de los componentes y se ha configurado el sistema completo. En primer lugar, se proporciona una descripción general de los sensores seleccionados y se explica la importancia de monitorear los contaminantes en el aire debido a sus efectos sobre la salud humana. En segundo lugar, se llevó a cabo un monitoreo en diferentes sitios de Bogotá, en particular en las sedes de la universidad Distrital Francisco José de Caldas, donde se instalaron los concentradores junto con sus antenas. La instalación de estos dispositivos requería de energía y conectividad, por lo que se realizó una evaluación detallada de los sitios de instalación y la gestión de permisos correspondientes. Se realizó un muestreo para determinar la cobertura de la red, y se presentaron gráficos y mapas que ilustran la distribución de los sensores en la ciudad.

Posteriormente, se llevaron a cabo análisis de los datos obtenidos por los sensores y se compararon con los valores mínimos de contaminantes. Finalmente, se comparó el modelo implementado con otros modelos ya certificados de fabricantes como el sistema \gls{lwan} Multitech e ISMT, y se realizaron pruebas de interoperabilidad que resultaron exitosas. Este capítulo concluye con una evaluación detallada de los resultados obtenidos y una discusión de las implicaciones de los hallazgos para la gestión de la calidad del aire en la ciudad, la cobertura y la interoperabilidad.

\chapter{Generalidades de la redes IoT}

A pesar de que el concepto de \gls{iot}, por sus siglas en inglés, es relativamente reciente, habiendo aparecido por primera vez entre los años 2008 y 2009 \cite{Evans2011TheEverything}, y definido por el \gls{ibsg} de Cisco como “el punto en el tiempo en el que se conectaron a internet más cosas u objetos que personas” \cite{Alcaraz2014InternetCosas}, en la actualidad el IoT tiene un gran impacto en la vida cotidiana. Cada vez se desarrollan más dispositivos con conexión a internet, de modo que no solo las personas se conectan a la red, sino que también los objetos necesitan estar conectados para funcionar correctamente o cumplir con sus propósitos

\subsection*{Sobre la evaluación de la red IoT propuesta}

El \gls{iot} tiene aplicaciones tanto medioambientales como industriales, así como en ciudades, en el hogar y a nivel personal. Las aplicaciones medioambientales incluyen sensores orientados a la protección y salud, no solo del ser humano, sino también de nuestro planeta. Un ejemplo es el monitoreo de la polución en mares y ríos, midiendo el estado del agua, su PH, salinidad, temperatura, iones, entre otros. Por otro lado, existen sistemas dedicados a la protección de animales mediante geolocalización y aplicaciones para la gestión de catástrofes. Los sensores que tradicionalmente solo podían ser leídos en sus centros de control, ahora pueden comunicarse por medio de GSM y ser fácilmente accesibles (entre ellos y con las centrales). Un ejemplo destacado es el proyecto ALARMS de la BGS (British Geological Survey, 2010), que ha logrado integrarse en un dispositivo de bajo costo con capacidad GSM y sensores miniaturizados para alertar tempranamente sobre la inestabilidad del suelo, anticipando así los deslizamientos de tierra \cite{Valle2014}.

En el ámbito industrial es donde se desarrollan más aplicaciones; algunas de las actuales incluyen la monitorización de estructuras, seguridad y análisis de datos (sistemas de big data), entre otras. Las ciudades también cuentan con aplicaciones tecnológicas basadas en IoT, como en los estacionamientos, contenedores inteligentes y el control eléctrico.

En el hogar, las aplicaciones se centran en gran medida en el control del entorno, como los sistemas de riego autónomo para plantas. Por su parte, las aplicaciones personales abarcan una amplia variedad de desarrollos, incluyendo píldoras inteligentes y trajes de monitoreo para bebés \cite{Valle2014}.

Estas nuevas tecnologías, puestas a nuestra disposición, pueden ser utilizadas en casi todos los aspectos de la vida cotidiana y en la industria, gracias a las numerosas posibilidades que ofrece el IoT para obtener, procesar y transmitir información. Cuando se destaca la versatilidad de estas tecnologías en distintos ámbitos de la vida diaria, se incluyen aplicaciones como el control de electrodomésticos o cualquier otro dispositivo en el hogar capaz de recoger y/o procesar datos, monitorear el estado de la vivienda e incluso verificar nuestra salud \cite{Alcaraz2014InternetCosas}. No obstante, las aplicaciones mencionadas anteriormente benefician principalmente a los propietarios, pero también pueden extenderse a la comunidad mediante soluciones de gran alcance como redes inteligentes, luces de la calle conectadas y autos inteligentes.

\section*{Sobre la alianza \gls{lwan}}

Los sistemas de comunicaciones nacieron como una solución a las necesidades de las personas para comunicarse de forma remota; luego de  su evolución, han trascendido a las máquinas y últimamente a los objetos de interés. La comunicación de objetos se inicio para extender las redes tipo \gls{lan} y sus derivados. Luego se vio la necesidad de comunicar las máquinas de forma remota para realizar procesos síncronos, y fue allí donde surgieron las redes de telecomunicaciones celulares para prestar este servicio, creándose protocolos estandarizados para tal fin, como el M2M. La conexión de objetos bajo este protocolo se hizo costosa y no tuvo acogida. Para bajar costos y poder conectar objetos de forma masiva, hubo iniciativas privadas como Sigfox, Kore, Weighless, etc. Estas compañías se basan en tecnologías de baja potencia de área amplia LP-WAN. El éxito fue inmediato pero solo eran iniciativas empresariales.

Fue así como se inicio una alianza de los principales sectores: fabricantes, operadores, usuarios; para crear una tecnología abierta tipo  LP-WAN, la cual se llamó \gls{lora}, donde cualquier persona o empresa pueda fabricar, implementar o prestar servicios para la conexión de objetos. Las características fundamentales de LoRa son: bajo consumo de potencia, tasa de bit baja, hasta 15 kilómetros de enlace, utilización de bandas libres ISM, entre otras. 

Como parte de los avances tecnológicos, en el área del internet de las cosas, emergen nuevas propuestas que mejoran las aplicaciones actuales, es el caso de las tecnologías \gls{lpwan} \cite{LoRa-2015}, las cuales complementan las tecnologías celulares y las de corto alcance para darle inteligencia a los objetos. Aun cuando los proveedores de telefonía celular también han iniciado el mercado en este segmento,  y en especial para la tecnología \gls{m2m}, la cual puede ser extendida a la conexión de objetos. Sin embargo, se necesitaron de varios años para lograr que la industria y los operadores se pusiera de acuerdo para crear los estándares que soportaran esta tecnología, para tal fin se creo la \textit{LoRa Alliance} para definir un estándar global para redes \gls{iot} con tecnología \gls{lpwan}, \gls{m2m} y para aplicaciones industriales o de consumo.

Este proyecto es el resultado de un proceso de investigación para las redes de internet de las cosas de área amplia con tecnología \gls{lpwan}, bajo el estándar \gls{lora}. El objeto del proyecto es evaluar la red en mención, puesto  que es un estándar nuevo y de unas prestaciones que mejoran los estándares que hay en el mercado, dado que es una alianza de cientos de empresas para la conexión de objetos con sistemas tecnológicos abiertos, para la elaboración, fabricación y prestación de servicios IoT sin restricciones.

Como punto de partida, se hizo un estudio de las técnicas utilizadas en la conformación de los estándares, y lo mas relevante es el uso eficiente que se hace de la energía dado que, según las tendencias en la conexión de objetos, el crecimiento es exponencial. Para ir en consonancia en el bajo consumo de energía, se incluyo la técnica de transmisión de datos  \gls{fhss}, por su eficiencia en el manejo del ruido, como la capa física. la capa \gls{mac} que maneja las políticas del canal y la capa de aplicación. 

La red se compone de una pasarela (gateway) a la cual se le pueden conectar los nodos que le dan inteligencia a los objetos. Para el caso de este proyecto se seleccionaron como objetos, un conjunto de sensores para monitorizar ciertas condiciones del ambiente y del aire.

El estándar LoRa se define como la capa física o la modulación inalámbrica utilizada para crear enlaces de comunicaciones de rango amplio. La técnica por desplazamiento de fase, \gls{fsk}, y modulado por ensanchamiento de espectro filtrado, \gls{css}, se utiliza para  la comunicación entre nodos y pasarelas. 

La comunicación entre los nodos y la pasarela son ensanchadas en diferentes canales y con distintas tasas de datos. La velocidad de los datos está entre los 0.3 y los 50 kbps, donde los nodos se pueden programar con distintas tasas de datos en una misma pasarela, esto es posible puesto que la red LoRa puede gestionar la velocidad de los datos y la salida para cada nodo de la red por medio de un esquema adaptativo de velocidad de datos, \gls{adr}.

Algunos parámetros se especifican por regulaciones regionales y por la clase de protocolo LoRa, el cual especifica las condiciones de los protocolos de acuerdo a su aplicación, como por ejemplo, si es \textit{duplex} o \textit{fullduplex}, o el tipo de velocidades de los datos.

Los tipos de comunicaciones LoRa se enmarcan en clase A, clase B y Clase C. actualmente esta terminada y en etapa de expansión la clase A, donde los dispositivos, pasarelas y nodos, están en producción. Las otras clases están en experimentación y para este documento se hará el trabajo con los dispositivos de la clase A.

\textbf{LoRa Clase A}: En esta clase, la comunicación es bidireccional, de tal forma que a cada nodo, cuando se comunica con el servidor (pasarela), se le permite la recepción de dos ventanas de tiempo. La transmisión de las ranuras  esta basado en las necesidades de cada nodo, donde pueden tener pequeñas variaciones aleatorias similares a los protocolos de transmisión ALOHA.   luego el nodo se coloca en modo escucha y después de un retardo corto espera en dos ventanas de tiempo para recibir comunicación del servido; si recibe información en la primera ventana el nodo termina el modo de escucha. Luego este debe esperar para trasmitir datos de acuerdo a una programación definida.

\textbf{LoRa Clase B}: Es una ampliación de la clase A, donde se permite programar las recepción de mas de dos ventanas por cada nodo. Esto permite abrir ventanas extras de tiempo para un nodo. La programación de recepción para estas ventanas extras se hace con una programación en la pasarela que permite la sincronización tipo Beacon, en la cual se envía una alerta por radiodifusión desde el servidor a los nodos que están en modo escucha. En general es una comunicación bidireccional entre la pasarela y los nodos con ranuras de tiempo de recepción programadas.

\textbf{LoRa Clase C}: Para esta clase, se tienen ventanas de recepción abiertas, casi continuamente y solo están cerradas cuando el nodo transmite. Esto contrasta con un mayor consumo de potencia con respecto a las clases anteriores.
